\section{Hand-to-Mouth Workers}\label{sec: hand_to_mouth_workers}

\q{HandToMouthWorkers}{
This section considers a version of the model in which labor is supplied by
hand-to-mouth workers rather than by investors.
We show that this modification leaves our main results unchanged.
In particular, in the special case $\nu = 0$ considered in
Section~\ref{sec:analysis.sec4}, the equilibrium allocation and asset prices
are identical to those in the baseline model.

\paragraph{Environment.}
The economy is populated by two types of households: investors and workers.

There is a finite number of investor types, indexed by $i \in \mathcal{I} = \{1,\ldots,I\}$,
with masses $\{\mu_i\}$ satisfying $\sum_i \mu_i = 1$.
Investors behave exactly as in the baseline model, except that we assume
their labor disutility is infinite, so that they supply no labor in equilibrium:
$h_{i,t} = 0$.

In addition, there is a unit mass of workers.
Workers are hand-to-mouth and consume all their labor income:
\[
C_{w,t} = W_t h_{w,t}.
\]
Workers have GHH preferences, and their labor supply condition is
\[
W_t = \xi_t h_{w,t}^{\nu},
\qquad
\xi_t = \xi A_{t-1},
\]
as in the baseline model.

\paragraph{Characterization.}
Investor $i$ solves Problem~\eqref{eq:recursive_representation}.
Hence, the characterization of investors' portfolio and consumption choices
is unchanged relative to the baseline model.
The firm's problem is also unchanged.

Labor market clearing implies
\[
h_t = h_{w,t}.
\]
Since the worker labor supply condition coincides with the labor supply
equation in the baseline model, the equilibrium determination of hours and wages
is identical.

Because workers do not save, financial market clearing conditions are unchanged:
\[
\sum_{i=1}^I \mu_i B_{i,t} = 0,
\qquad
\sum_{i=1}^I \mu_i S_{i,t} = 1.
\]

The only difference arises in goods market clearing:
\[
C_{w,t} + \sum_{i=1}^I \mu_i C_{i,t}
=
A_t h_t^{\alpha}.
\]
Using $C_{w,t} = W_t h_t = \xi_t h_t^{1+\nu}$,
we obtain
\[
\sum_{i=1}^I \mu_i C_{i,t}
=
A_t h_t^{\alpha} - \xi_t h_t^{1+\nu}.
\]

In the baseline model, where investors supply labor,
the goods market clearing condition can be written in terms of net consumption
$\tilde C_{i,t} \equiv C_{i,t} - \xi_t \frac{h_t^{1+\nu}}{1+\nu}$:
\[
\sum_{i=1}^I \mu_i \tilde C_{i,t}
=
A_t h_t^{\alpha}
-
\xi_t \frac{h_t^{1+\nu}}{1+\nu}.
\]

\paragraph{Equivalence when $\nu = 0$.}

When $\nu = 0$, labor income equals labor disutility in the baseline model.
In this case, the two goods market clearing conditions coincide exactly.
Hence, equilibrium allocations, asset prices, wealth dynamics,
and the evolution of market beliefs are identical to those in the baseline economy.

\paragraph{Economic intuition.}

Separating workers and investors therefore does not alter the mechanism
driving our results.
What matters is the timing of hiring, not who supplies labor.

When $\nu > 0$, investor consumption becomes more volatile in the
hand-to-mouth specification, because labor income accrues to workers,
whose consumption is smoother than output.
As a result, investors absorb a larger share of aggregate risk, consistent with the evidence on the consumption of stockholders (see, e.g., \citealt{mankiw1991consumption} and \citealt{parker2001consumption}).\footnote{This mechanism is reminiscent of, e.g., \citet{danthine2002labour} where limited market participation of workers, combined with firms insuring workers against labor risk, creates a similar operating-leverage channel.}
This \textit{strengthens} the operating-leverage channel without overturning
any of the qualitative results.

More generally, the mechanism linking beliefs, risk premia,
and labor demand operates through asset pricing and the firm's first-order condition,
and does not rely on investors supplying labor directly.}

